\subsection{$[[5,1,2]]_3$ Qutrit Surface Code}

\subsubsection{20-Qutrit IQM Garnet QPU}

Our experiments are performed on IQM's superconducting 20-qutrit Garnet QPU.  Each physical unit is a transmon circuit operated as a three-level system with computational basis states $\{\ket{0},\ket{1},\ket{2}\}$, coupled via tunable superconducting couplers that generate a ZZ-type interaction.  By applying a magnetic
flux pulse to the coupler one implements a fast, resonant CZ-type gate with typical durations in the range of $20$--$40\,\mathrm{ns}$, while
the effective ZZ interaction can be turned off during idle periods. Single-qutrit Clifford gates are implemented by microwave drives on the
% TODO(authors): parametric drive or baseband flux pulse for the Garnet CZ? The previous sentence and the appendix sec:baseband-gates describe a flux pulse on the coupler; that appendix is not referenced from the main text.
$0\!\leftrightarrow\!1$ and $1\!\leftrightarrow\!2$ transitions, and two-qutrit entangling gates are realized using parametrically driven controlled-phase operations on the couplers. \textcolor{red}{ASK NOAH FOR FIGURES AND CONTENT REGARDING IQM QPU}.

From this device we select a connected $5$-qutrit patch and encode a single logical qutrit into these five data qutrits using the $[[5,1,2]]_3$ error-detection code defined in \autoref{sec: error detection code}. We then implement transversal Clifford gates within a modified LRB protocol to garner estimates for the average error rate of our logical Clifford gates. However, we don't use the standard qudit surface code and folded code definitions for stabilizers and logical operators, as shown in \autoref{sec: qudit surface code}, \cite{Moussa2016} and \autoref{sec: error detection code}. In \autoref{sec: error detection code} we explain our modifications and the motivation behind them. 


\subsubsection{$[[5,1,2]]_3$ Qutrit Error Detection Code} \label{sec: error detection code}

The logical qutrit is encoded using the $[[5,1,2]]_3$ quantum error-detection code obtained from a surface-code patch with five data qutrits and four stabilizer generators \cite{Moussa2016}. Let
$X_j$ and $Z_j$ denote the generalized Pauli operators acting on
qutrit $j$, with $X_j^{-1}=X_j^\dagger$ and $Z_j^{-1}=Z_j^\dagger$. The
code space is defined as the simultaneous $+1$ eigenspace of
\begin{align}
  G_1 &= Z_0 \otimes Z_1^{-1} \otimes Z_3^{-1}, \label{eq:s1}\\
  G_2 &= Z_1 \otimes Z_2 \otimes Z_4^{-1}, \\
  G_3 &= X_0 \otimes X_1 \otimes X_2^{-1}, \\
  G_4 &= X_1^{-1} \otimes X_3 \otimes X_4^{-1},
  \label{eq:s4}
\end{align}
yielding a $[[5,1,2]]_3$ stabilizer code: five physical qutrits, one logical qutrit, and distance $d=2$. A convenient choice of logical Pauli operators is
\begin{equation}
  \bar Z = Z_0 \otimes Z_2, \qquad
  \bar X = X_0 \otimes X_3,
  \label{eq:logicalPauli}
\end{equation}
which commute with all $G_i$ and generate the logical Pauli group modulo the stabilizer.

Following the folded-surface-code construction of
\cite{Moussa2016}, we realize a logical single-qutrit Clifford generating set $\{\bar H,\bar S\}$ in constant-depth transversal form. In our hardware-aware
implementation we use
\begin{align}
  \bar H &=
  (H_0 \otimes H_1^{-1} \otimes H_2 \otimes H_3 \otimes H_4)
  (I_0 \otimes I_1 \otimes \mathrm{SWAP}_{2,3} \otimes I_4), \\
  \bar S &= S_0 \otimes S_1^{-1} \otimes CZ_{2\rightarrow 3} \otimes S_4,
  \label{eq:logical_HS}
\end{align}
% TODO(authors): SWAP_{2,3} and CZ_{2->3} couple q2 and q3. In the q2-centered patch of fig: 5-qutrit-code and in the two-ancilla layout (fig:ancilla_layouts(c)) they are diagonal neighbors, but in the one-ancilla layout (fig:ancilla_layouts(b), q1 at the center) they sit at opposite corners and need a long-range coupler. Please check that "we avoid any long range 2-qutrit gates" below, and "nearest-neighbor" in the caption of fig: 5-qutrit-code, still read correctly.
where subscripts label physical qutrits, $H$ and $S$ are single-qutrit Clifford generators, and $CZ_{2\rightarrow 3}$ is a controlled-phase gate with control on qutrit~2 and target on qutrit~3. Compared to the canonical folded-code realization, these definitions are adjusted to minimize the number of local $\mathrm{SWAP}$ operations required on the QPU. As shown in \autoref{fig: logical H} and \autoref{fig: logical S}, on the patch of \autoref{fig: 5-qutrit-code}, with $q_2$ at the center, we avoid any long range 2-qutrit gates when defining the logical Clifford gates and therefore avoid the folding operation altogether, described in \autoref{sec: qudit surface code} and \cite{Moussa2016}. The conjugation action of the logical Cliffords give us the following conjugation relations: writing $G_Z$ ($G_X$) for a generic $Z$-type ($X$-type) stabilizer generator, we have
\begin{table}[H]
  \centering
  \begin{tabular}{c|c}
    \hline
    Conjugation of Stabilizers
      & Conjugation of Logical Paulis\\
    \hline
    $\bar H G_Z \bar H^\dagger = G_X^\dagger$
      & $\bar H \bar Z \bar H^\dagger = \bar X^\dagger$ \\
    $\bar H G_X \bar H^\dagger = G_Z$
      & $\bar H \bar X \bar H^\dagger = \bar Z$ \\
    $\bar S G_Z \bar S^\dagger = G_Z$
      & $\bar S \bar Z \bar S^\dagger = \bar Z$ \\
    $\bar S G_X \bar S^\dagger = G_X G_Z$
      & $\bar S \bar X \bar S^\dagger = \bar X \bar Z$ \\
    \hline
  \end{tabular}
  \caption{Conjugation action of the logical single qudit Clifford group generators on
  stabilizers (left) and logical Pauli operators (right) for the  $[[5,1,2]]_3$ code.}
  \label{tab:clifford-conjugation}
\end{table}
These relations show that $\bar H$ and $\bar S$ map the stabilizer group to itself and implement the usual Clifford action on the logical Pauli operators for the $[[5,1,2]]_3$ code.

% ============================================================
% fig - 5-qutrit code.tex
% Figure: [[5,1,2]]_3 patch definitions, logical Cliffords, stabilizers,
% and logical Paulis (d=2 nearest-neighbour construction).
% All macros come from 'tikz preamble.tex'.
% ============================================================

\begin{figure}[H]
\centering
% ===================== (a) Reference =====================
\begin{subfigure}[b]{0.9\textwidth}
\centering
\begin{tikzpicture}[line cap=round,line join=round,thick,scale=1.0]
  \coordinate (Left)   at (-4.85, 0);
  \coordinate (Middle) at (-1.85, 0);
  \coordinate (XGrid)  at ( 1.55, 0);
  \coordinate (ZGrid)  at ( 4.45, 0);

  \node[above,font=\small] at ($(Left)+(0,1.72)$) {Unfolded};
  \node[above,font=\small] at ($(Middle)+(0,1.72)$) {Folded};
  \node[above,font=\small] at ($(XGrid)!0.5!(ZGrid)+(0,1.80)$) {Redefined Stabilizers};
  \node[above,font=\footnotesize,text=Xcol!80!black] at ($(XGrid)+(0,1.20)$) {$X$-type Stabilizers};
  \node[above,font=\footnotesize,text=Zcol!85!black] at ($(ZGrid)+(0,1.20)$) {$Z$-type Stabilizers};
  \node[font=\scriptsize,text=Xcol!80!black] at ($(XGrid)+(0,-1.25)$) {$G_3,\;G_4$};
  \node[font=\scriptsize,text=Zcol!85!black] at ($(ZGrid)+(0,-1.25)$) {$G_1,\;G_2$};

  % --- Left Patch (Unfolded) ---
  \begin{scope}[shift={(Left)}]
    \FiveQPatchFilled
    \FiveQPatchCircles
  \end{scope}

  % --- Middle Patch (Folded) ---
  \begin{scope}[shift={(Middle)}]
    \FiveQPatchOverlapped
    \def\m{0.30}
    \draw[dashed, draw=red, line width=1pt, dash pattern=on 3pt off 2pt]
      ($(q0)+(-\m,\m)$) -- ($(q4)+(\m,-\m)$);
    \FiveQPatchCirclesOneUnderThree
  \end{scope}

  \draw[-{Stealth[length=3mm,width=2mm]},red,line width=1.6pt]
    ($(Middle)+(1.35,0)$) -- ($(XGrid)+(-1.35,0)$);

  % --- Redefined X-type grid ---
  \begin{scope}[shift={(XGrid)}]
    \FiveQPatchRedefinedXGrid
    \FiveQPatchCircles
  \end{scope}

  % --- Redefined Z-type grid ---
  \begin{scope}[shift={(ZGrid)}]
    \FiveQPatchRedefinedZGrid
    \FiveQPatchCircles
  \end{scope}
\end{tikzpicture}
\caption{5-Qutrit Surface Code Patch}
\end{subfigure}

\vspace{2.0ex}

% ===================== (b) Definition =====================
\hspace*{-1.5em}
\begin{subfigure}[c]{0.48\textwidth}
\centering
\begin{tikzpicture}[line cap=round,line join=round,thick,scale=1.0]
  \coordinate (LeftTop) at (-2.2,  1.50);
  \coordinate (LeftBot) at (-2.2, -1.50);
  \coordinate (RightTop) at ( 2.2,  1.50);
  \coordinate (RightBot) at ( 2.2, -1.50);

  \node[scale=1.0, anchor=east] at ($(LeftTop)+(-1.4, -1.5)$) {$\bar{H'}\!:=$};
  \node[scale=1.0, anchor=east] (HLabel) at ($(RightTop)+(-1.3, -1.5)$) {$\bar{H}\!:=$};

  \def\lw{1.5pt} \def\rc{8pt}
  % LEFT BOX 
  \path let \p1 = ($(LeftTop)+(-1.3,1.7)$), \p2 = ($(LeftBot)+(1.4,-1.4)$)
        in coordinate (LTL) at (\x1,\y1) coordinate (LBR) at (\x2,\y2)
           coordinate (LTR) at (\x2,\y1) coordinate (LBL) at (\x1,\y2);
  \coordinate (LMidLeft) at ($(LTL)!0.5!(LBL)$); \coordinate (LMidRight) at ($(LTR)!0.5!(LBR)$);
  \draw[HCol, dash pattern=on 2pt off 5pt, line width=\lw, rounded corners=\rc] (LMidLeft)--(LTL)--(LTR)--(LMidRight);
  \draw[SwapCol, dash pattern=on 2pt off 5pt, line width=\lw, rounded corners=\rc] (LMidRight)--(LBR)--(LBL)--(LMidLeft);

  % RIGHT BOX
  \path let \p1 = ($(RightTop)+(-1.3,1.7)$), \p2 = ($(RightBot)+(1.4,-1.4)$)
        in coordinate (RTL) at (\x1,\y1) coordinate (RBR) at (\x2,\y2)
           coordinate (RTR) at (\x2,\y1) coordinate (RBL) at (\x1,\y2);
  \coordinate (RMidLeft) at ($(RTL)!0.5!(RBL)$); \coordinate (RMidRight) at ($(RTR)!0.5!(RBR)$);
  \draw[HCol, dash pattern=on 2pt off 5pt, line width=\lw, rounded corners=\rc] (RMidLeft)--(RTL)--(RTR)--(RMidRight);
  \draw[SwapCol, dash pattern=on 2pt off 5pt, line width=\lw, rounded corners=\rc] (RMidRight)--(RBR)--(RBL)--(RMidLeft);

  \begin{scope}[shift={(LeftTop)}]
    \FiveQPatchEdges
    \foreach \P/\txt in {q0/0,q1/1,q2/2,q3/3,q4/4}{ \draw[line width=\qdLine,fill=HCol] (\P) circle (\qdR); \node[font=\numFont, text=white] at (\P) {\txt}; }
    \WedgeArrow{q0}{0}\WedgeArrow{q1}{0}\WedgeArrow{q3}{0}\WedgeArrow{q4}{0}\WedgeArrow{q2}{90}
    \node[anchor=south, font=\small, yshift=0mm] at (0, 1.0) {Layer 1};
    \begin{scope}[text=HCol, font=\footnotesize]
       \node[below, yshift=-3pt] at (q0) {$H_0$}; \node[below, yshift=-3pt] at (q1) {$H_1$};
       \node[right, xshift=2pt]  at (q2) {$H_2^{-1}$}; \node[above, yshift=3pt]  at (q3) {$H_3$}; \node[above, yshift=3pt]  at (q4) {$H_4$};
    \end{scope}
  \end{scope}

  \begin{scope}[shift={(LeftBot)}]
    \FiveQPatchEdges
    \foreach \P/\txt in {q0/0,q2/2,q4/4}{ \draw[line width=\qdLine,fill=white] (\P) circle (\qdR); \node[font=\numFont] at (\P) {\txt}; }
    \foreach \P/\txt in {q1/1,q3/3}{ \draw[line width=\qdLine,fill=SwapCol] (\P) circle (\qdR); \node[font=\numFont, text=white] at (\P) {\txt}; }
    \WedgeArrow{q0}{0}\WedgeArrow{q1}{0}\WedgeArrow{q3}{0}\WedgeArrow{q4}{0}\WedgeArrow{q2}{90}
    \node[anchor=south, font=\small, yshift=0mm] at (0, 1.0) {Layer 2};
    \draw[{Stealth[length=3mm,width=2mm]}-{Stealth[length=3mm,width=2mm]}, color=SwapCol, line width=1.2pt, shorten >=5pt, shorten <=3pt]
      (q1) to[bend left=25] node[midway, below right, rotate=52, xshift=-7pt, font=\small] {SWAP} (q3);
  \end{scope}

  \draw[-{Stealth[length=3mm,width=2mm]},red!75,line width=1.6pt] (-0.6, 0) -- (-0.1, 0);

  \begin{scope}[shift={(RightTop)}]
    \FiveQPatchEdges
    \foreach \P/\txt in {q0/0,q1/1,q2/2,q3/3,q4/4}{ \draw[line width=\qdLine,fill=HCol] (\P) circle (\qdR); \node[font=\numFont, text=white] at (\P) {\txt}; }
    \WedgeArrow{q0}{0}\WedgeArrow{q1}{0}\WedgeArrow{q3}{0}\WedgeArrow{q4}{0}\WedgeArrow{q2}{90}
    \node[anchor=south, font=\small, yshift=0mm] at (0, 1.0) {Layer 1};
    \begin{scope}[text=HCol, font=\footnotesize]
       \node[below, yshift=-3pt] at (q0) {$H_0$}; \node[below, yshift=-3pt] at (q1) {$H_1^{-1}$};
       \node[right, xshift=2pt]  at (q2) {$H_2$}; \node[above, yshift=3pt]  at (q3) {$H_3$}; \node[above, yshift=3pt]  at (q4) {$H_4$};
    \end{scope}
  \end{scope}

  \begin{scope}[shift={(RightBot)}]
    \FiveQPatchEdges
    \foreach \P/\txt in {q0/0,q1/1,q4/4}{ \draw[line width=\qdLine,fill=white] (\P) circle (\qdR); \node[font=\numFont] at (\P) {\txt}; }
    \foreach \P/\txt in {q2/2,q3/3}{ \draw[line width=\qdLine,fill=SwapCol] (\P) circle (\qdR); \node[font=\numFont, text=white] at (\P) {\txt}; }
    \WedgeArrow{q0}{0}\WedgeArrow{q1}{0}\WedgeArrow{q3}{0}\WedgeArrow{q4}{0}\WedgeArrow{q2}{90}
    \node[anchor=south, font=\small, yshift=0mm] at (0, 1.0) {Layer 2};
    \draw[{Stealth[length=3mm,width=2mm]}-{Stealth[length=3mm,width=2mm]}, color=SwapCol, line width=1.2pt, shorten >=3pt, shorten <=3pt]
      (q2) to[bend right=45] node[midway, below, xshift =-7pt, yshift=14pt, rotate=41, font=\small] {SWAP} (q3);
  \end{scope}
\end{tikzpicture}
\caption{Logical $H$}
\label{fig: logical H}
\end{subfigure}
\quad
% ===================== (c) S Definition =====================
\begin{subfigure}[c]{0.48\textwidth}
\centering
\begin{tikzpicture}[line cap=round,line join=round,thick,scale=1.0]
  \coordinate (Left)  at (-2.2, 0);
  \coordinate (Right) at ( 2.2, 0);
  \def\lw{1.5pt} \def\rc{8pt}

  \coordinate (LTL) at ($(Left)+(-1.3, 1.4)$); \coordinate (LBR) at ($(Left)+(1.4, -1.4)$);
  \coordinate (LTR) at ($(Left)+(1.4, 1.4)$);  \coordinate (LBL) at ($(Left)+(-1.3, -1.4)$);
  \coordinate (LMidLeft) at ($(LTL)!0.5!(LBL)$); \coordinate (LMidRight) at ($(LTR)!0.5!(LBR)$);
  \draw[HCol, dash pattern=on 2pt off 5pt, line width=\lw, rounded corners=\rc] (LMidLeft)--(LTL)--(LTR)--(LMidRight);
  \draw[SwapCol, dash pattern=on 2pt off 5pt, line width=\lw, rounded corners=\rc] (LMidRight)--(LBR)--(LBL)--(LMidLeft);

  \coordinate (RTL) at ($(Right)+(-1.3, 1.4)$); \coordinate (RBR) at ($(Right)+(1.4, -1.4)$);
  \coordinate (RTR) at ($(Right)+(1.4, 1.4)$);  \coordinate (RBL) at ($(Right)+(-1.3, -1.4)$);
  \coordinate (RMidLeft) at ($(RTL)!0.5!(RBL)$); \coordinate (RMidRight) at ($(RTR)!0.5!(RBR)$);
  \draw[HCol, dash pattern=on 2pt off 5pt, line width=\lw, rounded corners=\rc] (RMidLeft)--(RTL)--(RTR)--(RMidRight);
  \draw[SwapCol, dash pattern=on 2pt off 5pt, line width=\lw, rounded corners=\rc] (RMidRight)--(RBR)--(RBL)--(RMidLeft);

  \node[scale=1.0, anchor=east] at ($(Left)+(-1.4, 0)$) {$\bar{S'}\!:=$};
  \node[scale=1.0, anchor=east] (SLabel) at ($(Right)+(-1.4, 0)$) {$\bar{S}\!:=$};

  \newcommand{\DrawLHSPattern}{
      \FiveQPatchEdges
      \foreach \P in {q0, q2, q4} { \draw[line width=\qdLine,fill=HCol] (\P) circle (\qdR); }
      \foreach \P in {q1, q3} { \draw[line width=\qdLine,fill=SwapCol] (\P) circle (\qdR); }
      \foreach \P/\txt in {q0/0,q1/1,q2/2,q3/3,q4/4}{ \node[font=\numFont, text=white] at (\P) {\txt}; }
      \WedgeArrow{q0}{0}\WedgeArrow{q1}{0}\WedgeArrow{q3}{0}\WedgeArrow{q4}{0}\WedgeArrow{q2}{90}
      \node[text=HCol, font=\small, below, yshift=-3pt] at (q0) {$S_0$};
      \node[text=HCol, font=\small, right, xshift=3pt] at (q2) {$S_2^{-1}$};
      \node[text=HCol, font=\small, above, yshift=3pt] at (q4) {$S_4$};
      \draw[color=SwapCol, line width=1.2pt, shorten <=4pt, shorten >=4pt] 
         (q1) to[out=180, in=90] coordinate[pos=0.45] (mid) (q3);
      \node[text=SwapCol, font=\small, rotate=52, anchor=south, xshift=-12pt, yshift=-1pt] at (mid) {$CZ_{1\to 3}$};
  }
  \newcommand{\DrawRHSPattern}{
      \FiveQPatchEdges
      \foreach \P in {q0, q1, q4} { \draw[line width=\qdLine,fill=HCol] (\P) circle (\qdR); }
      \foreach \P in {q2, q3} { \draw[line width=\qdLine,fill=SwapCol] (\P) circle (\qdR); }
      \foreach \P/\txt in {q0/0,q1/1,q2/2,q3/3,q4/4}{ \node[font=\numFont, text=white] at (\P) {\txt}; }
      \WedgeArrow{q0}{0}\WedgeArrow{q1}{0}\WedgeArrow{q3}{0}\WedgeArrow{q4}{0}\WedgeArrow{q2}{90}
      \node[text=HCol, font=\small, below, yshift=-3pt] at (q0) {$S_0$};
      \node[text=HCol, font=\small, below, yshift=-3pt] at (q1) {$S_1^{-1}$};
      \node[text=HCol, font=\small, above, yshift=3pt] at (q4) {$S_4$};
      \draw[color=SwapCol, line width=1.2pt, shorten <=4pt, shorten >=4pt] 
         (q3) to[out=90, in=180] coordinate[pos=0.5] (midR) (q2);
      \node[text=SwapCol, font=\small, rotate=40, anchor=south, xshift=5pt, yshift=-2pt] at (midR) {$CZ_{2\to 3}$};
  }
  \begin{scope}[shift={(Left)}] \DrawLHSPattern \end{scope}
  \begin{scope}[shift={(Right)}] \DrawRHSPattern \end{scope}
  \draw[-{Stealth[length=3mm,width=2mm]},red,line width=1.6pt] (-0.6, 0) -- (-0.1, 0);
\end{tikzpicture}
\caption{Logical $S$}
\label{fig: logical S}
\end{subfigure}

\vspace{2.0ex}

% ===================== (d) Parallel Operations =====================
\begin{subfigure}[b]{0.9\textwidth}
\centering
\begin{tikzpicture}[line cap=round,line join=round,thick,scale=1.0]
  \def\yTop{1.6} \def\yBot{-1.2}

  % COL 1
  \def\x{-4.8}
  \begin{scope}[shift={(\x, \yTop)}]
     \FiveQPatchEdges
     \foreach \P/\txt in {q1/1, q4/4}{ \draw[line width=\qdLine,fill=white] (\P) circle (\qdR); \node[font=\numFont] at (\P) {\txt}; }
     \foreach \P/\txt in {q0/0, q2/2, q3/3}{ \draw[line width=\qdLine,fill=Zcol] (\P) circle (\qdR); \node[font=\numFont, text=white] at (\P) {\txt}; }
     \WedgeArrow{q0}{0}\WedgeArrow{q1}{0}\WedgeArrow{q3}{0}\WedgeArrow{q4}{0}\WedgeArrow{q2}{90}
     \begin{scope}[text=Zcol!85!black, font=\footnotesize]
        \node[below, yshift=-3pt] at (q0) {$Z_0$}; \node[right, xshift=2pt] at (q2) {$Z_2^{-1}$}; \node[above, yshift=3pt] at (q3) {$Z_3^{-1}$};
     \end{scope}
  \end{scope}
  \begin{scope}[shift={(\x, \yBot)}]
     \FiveQPatchEdges
     \foreach \P/\txt in {q2/2, q4/4}{ \draw[line width=\qdLine,fill=white] (\P) circle (\qdR); \node[font=\numFont] at (\P) {\txt}; }
     \foreach \P/\txt in {q0/0, q1/1, q3/3}{ \draw[line width=\qdLine,fill=Zcol] (\P) circle (\qdR); \node[font=\numFont, text=white] at (\P) {\txt}; }
     \WedgeArrow{q0}{0}\WedgeArrow{q1}{0}\WedgeArrow{q3}{0}\WedgeArrow{q4}{0}\WedgeArrow{q2}{90}
     \begin{scope}[text=Zcol!85!black, font=\footnotesize]
        \node[below, yshift=-3pt] at (q0) {$Z_0$}; \node[below, yshift=-3pt] at (q1) {$Z_1^{-1}$}; \node[above, yshift=3pt] at (q3) {$Z_3^{-1}$};
     \end{scope}
  \end{scope}
  \draw[-{Stealth[length=3mm,width=2mm]},red!75,line width=1.6pt] (\x, 0.425) -- (\x, -0.1);

  % COL 2
  \def\x{-1.6}
  \begin{scope}[shift={(\x, \yTop)}]
     \FiveQPatchEdges
     \foreach \P/\txt in {q0/0, q3/3}{ \draw[line width=\qdLine,fill=white] (\P) circle (\qdR); \node[font=\numFont] at (\P) {\txt}; }
     \foreach \P/\txt in {q1/1, q2/2, q4/4}{ \draw[line width=\qdLine,fill=Zcol] (\P) circle (\qdR); \node[font=\numFont, text=white] at (\P) {\txt}; }
     \WedgeArrow{q0}{0}\WedgeArrow{q1}{0}\WedgeArrow{q3}{0}\WedgeArrow{q4}{0}\WedgeArrow{q2}{90}
     \begin{scope}[text=Zcol!85!black, font=\footnotesize]
        \node[below, yshift=-3pt] at (q1) {$Z_1$}; \node[right, xshift=2pt] at (q2) {$Z_2$}; \node[above, yshift=3pt] at (q4) {$Z_4^{-1}$};
     \end{scope}
  \end{scope}
  \begin{scope}[shift={(\x, \yBot)}]
     \FiveQPatchEdges
     \foreach \P/\txt in {q0/0, q3/3}{ \draw[line width=\qdLine,fill=white] (\P) circle (\qdR); \node[font=\numFont] at (\P) {\txt}; }
     \foreach \P/\txt in {q1/1, q2/2, q4/4}{ \draw[line width=\qdLine,fill=Zcol] (\P) circle (\qdR); \node[font=\numFont, text=white] at (\P) {\txt}; }
     \WedgeArrow{q0}{0}\WedgeArrow{q1}{0}\WedgeArrow{q3}{0}\WedgeArrow{q4}{0}\WedgeArrow{q2}{90}
     \begin{scope}[text=Zcol!85!black, font=\footnotesize]
        \node[below, yshift=-3pt] at (q1) {$Z_1$}; \node[right, xshift=2pt] at (q2) {$Z_2$}; \node[above, yshift=3pt] at (q4) {$Z_4^{-1}$};
     \end{scope}
  \end{scope}
  \draw[-{Stealth[length=3mm,width=2mm]},red!75,line width=1.6pt] (\x, 0.425) -- (\x, -0.1);

  % COL 3
  \def\x{1.6}
  \begin{scope}[shift={(\x, \yTop)}]
     \FiveQPatchEdges
     \foreach \P/\txt in {q3/3, q4/4}{ \draw[line width=\qdLine,fill=white] (\P) circle (\qdR); \node[font=\numFont] at (\P) {\txt}; }
     \foreach \P/\txt in {q0/0, q1/1, q2/2}{ \draw[line width=\qdLine,fill=Xcol] (\P) circle (\qdR); \node[font=\numFont, text=white] at (\P) {\txt}; }
     \WedgeArrow{q0}{0}\WedgeArrow{q1}{0}\WedgeArrow{q3}{0}\WedgeArrow{q4}{0}\WedgeArrow{q2}{90}
     \begin{scope}[text=Xcol!80!black, font=\footnotesize]
        \node[below, yshift=-3pt] at (q0) {$X_0$}; \node[below, yshift=-3pt] at (q1) {$X_1^{-1}$}; \node[right, xshift=2pt] at (q2) {$X_2$};
     \end{scope}
  \end{scope}
  \begin{scope}[shift={(\x, \yBot)}]
     \FiveQPatchEdges
     \foreach \P/\txt in {q3/3, q4/4}{ \draw[line width=\qdLine,fill=white] (\P) circle (\qdR); \node[font=\numFont] at (\P) {\txt}; }
     \foreach \P/\txt in {q0/0, q1/1, q2/2}{ \draw[line width=\qdLine,fill=Xcol] (\P) circle (\qdR); \node[font=\numFont, text=white] at (\P) {\txt}; }
     \WedgeArrow{q0}{0}\WedgeArrow{q1}{0}\WedgeArrow{q3}{0}\WedgeArrow{q4}{0}\WedgeArrow{q2}{90}
     \begin{scope}[text=Xcol!80!black, font=\footnotesize]
        \node[below, yshift=-3pt] at (q0) {$X_0$}; \node[below, yshift=-3pt] at (q1) {$X_1$}; \node[right, xshift=2pt] at (q2) {$X_2^{-1}$};
     \end{scope}
  \end{scope}
  \draw[-{Stealth[length=3mm,width=2mm]},red!75,line width=1.6pt] (\x, 0.425) -- (\x, -0.1);

  % COL 4
  \def\x{4.8}
  \begin{scope}[shift={(\x, \yTop)}]
     \FiveQPatchEdges
     \foreach \P/\txt in {q0/0, q1/1}{ \draw[line width=\qdLine,fill=white] (\P) circle (\qdR); \node[font=\numFont] at (\P) {\txt}; }
     \foreach \P/\txt in {q2/2, q3/3, q4/4}{ \draw[line width=\qdLine,fill=Xcol] (\P) circle (\qdR); \node[font=\numFont, text=white] at (\P) {\txt}; }
     \WedgeArrow{q0}{0}\WedgeArrow{q1}{0}\WedgeArrow{q3}{0}\WedgeArrow{q4}{0}\WedgeArrow{q2}{90}
     \begin{scope}[text=Xcol!80!black, font=\footnotesize]
        \node[right, xshift=2pt] at (q2) {$X_2^{-1}$}; \node[above, yshift=3pt] at (q3) {$X_3$}; \node[above, yshift=3pt] at (q4) {$X_4^{-1}$};
     \end{scope}
  \end{scope}
  \begin{scope}[shift={(\x, \yBot)}]
     \FiveQPatchEdges
     \foreach \P/\txt in {q0/0, q2/2}{ \draw[line width=\qdLine,fill=white] (\P) circle (\qdR); \node[font=\numFont] at (\P) {\txt}; }
     \foreach \P/\txt in {q1/1, q3/3, q4/4}{ \draw[line width=\qdLine,fill=Xcol] (\P) circle (\qdR); \node[font=\numFont, text=white] at (\P) {\txt}; }
     \WedgeArrow{q0}{0}\WedgeArrow{q1}{0}\WedgeArrow{q3}{0}\WedgeArrow{q4}{0}\WedgeArrow{q2}{90}
     \begin{scope}[text=Xcol!80!black, font=\footnotesize]
        \node[below, yshift=-3pt] at (q1) {$X_1^{-1}$}; \node[above, yshift=3pt] at (q3) {$X_3$}; \node[above, yshift=3pt] at (q4) {$X_4^{-1}$};
     \end{scope}
  \end{scope}
  \draw[-{Stealth[length=3mm,width=2mm]},red!75,line width=1.6pt] (\x, 0.425) -- (\x, -0.1);
\end{tikzpicture}
\caption{Stabilizers for the $[[5,1,2]]_3$ code}
\end{subfigure}

\vspace{2.0ex}
% ===================== (e) Logical X =====================
\begin{subfigure}[c]{0.35\textwidth}
\centering
\begin{tikzpicture}[line cap=round,line join=round,thick,scale=1.0]
  \node[scale=1.0, anchor=east] at (-1.4, 0) {$\bar{X'}=\bar{X}\!:=$};
   
  % Added blue box
  \def\lw{1.5pt} \def\rc{8pt}
  \coordinate (LTL_X) at (-1.3, 1.4); \coordinate (LBR_X) at (1.4, -1.4);
  \coordinate (LTR_X) at ( 1.4, 1.4); \coordinate (LBL_X) at (-1.3, -1.4);
  \draw[Xcol, dash pattern=on 2pt off 5pt, line width=\lw, rounded corners=\rc] (LBL_X)--(LTL_X)--(LTR_X)--(LBR_X)--cycle;
   
  \FiveQPatchEdges
  \foreach \P/\txt in {q1/1, q2/2, q4/4}{ \draw[line width=\qdLine,fill=white] (\P) circle (\qdR); \node[font=\numFont] at (\P) {\txt}; }
  \foreach \P/\txt in {q0/0, q3/3}{ \draw[line width=\qdLine,fill=Xcol] (\P) circle (\qdR); \node[font=\numFont, text=white] at (\P) {\txt}; }
  \WedgeArrow{q0}{0}\WedgeArrow{q1}{0}\WedgeArrow{q3}{0}\WedgeArrow{q4}{0}\WedgeArrow{q2}{90}
  \begin{scope}[text=Xcol!80!black, font=\footnotesize]
     \node[below, yshift=-3pt] at (q0) {$X_0$};
     \node[above, yshift=3pt]  at (q3) {$X_3$};
  \end{scope}
\end{tikzpicture}
\caption{Logical $X$}
\end{subfigure}
\quad
% ===================== (f) Logical Z =====================
\begin{subfigure}[c]{0.6\textwidth}
\centering
\begin{tikzpicture}[line cap=round,line join=round,thick,scale=1.0]
  % Define centers
  \coordinate (Left)  at (-2.2, 0);
  \coordinate (Right) at ( 2.2, 0);

  % Labels
  \node[scale=1.0, anchor=east] at ($(Left)+(-1.4, 0)$) {$\bar{Z'}\!:=$};
  \node[scale=1.0, anchor=east] at ($(Right)+(-1.4, 0)$) {$\bar{Z}\!:=$};
   
  % Added red boxes
  \def\lw{1.5pt} \def\rc{8pt}
  % Left Box
  \coordinate (LTL_Z) at ($(Left)+(-1.3, 1.4)$); \coordinate (LBR_Z) at ($(Left)+(1.4, -1.4)$);
  \coordinate (LTR_Z) at ($(Left)+(1.4, 1.4)$); \coordinate (LBL_Z) at ($(Left)+(-1.3, -1.4)$);
  \draw[Zcol, dash pattern=on 2pt off 5pt, line width=\lw, rounded corners=\rc] (LBL_Z)--(LTL_Z)--(LTR_Z)--(LBR_Z)--cycle;

  % Right Box
  \coordinate (RTL_Z) at ($(Right)+(-1.3, 1.4)$); \coordinate (RBR_Z) at ($(Right)+(1.4, -1.4)$);
  \coordinate (RTR_Z) at ($(Right)+(1.4, 1.4)$); \coordinate (RBL_Z) at ($(Right)+(-1.3, -1.4)$);
  \draw[Zcol, dash pattern=on 2pt off 5pt, line width=\lw, rounded corners=\rc] (RBL_Z)--(RTL_Z)--(RTR_Z)--(RBR_Z)--cycle;

  % Left Patch: Z0 Z1
  \begin{scope}[shift={(Left)}]
     \FiveQPatchEdges
     \foreach \P/\txt in {q2/2, q3/3, q4/4}{ \draw[line width=\qdLine,fill=white] (\P) circle (\qdR); \node[font=\numFont] at (\P) {\txt}; }
     \foreach \P/\txt in {q0/0, q1/1}{ \draw[line width=\qdLine,fill=Zcol] (\P) circle (\qdR); \node[font=\numFont, text=white] at (\P) {\txt}; }
     \WedgeArrow{q0}{0}\WedgeArrow{q1}{0}\WedgeArrow{q3}{0}\WedgeArrow{q4}{0}\WedgeArrow{q2}{90}
     \begin{scope}[text=Zcol!85!black, font=\footnotesize]
        \node[below, yshift=-3pt] at (q0) {$Z_0$};
        \node[below, yshift=-3pt] at (q1) {$Z_1$};
     \end{scope}
  \end{scope}

  % Right Patch: Z0 Z2
  \begin{scope}[shift={(Right)}]
     \FiveQPatchEdges
     \foreach \P/\txt in {q1/1, q3/3, q4/4}{ \draw[line width=\qdLine,fill=white] (\P) circle (\qdR); \node[font=\numFont] at (\P) {\txt}; }
     \foreach \P/\txt in {q0/0, q2/2}{ \draw[line width=\qdLine,fill=Zcol] (\P) circle (\qdR); \node[font=\numFont, text=white] at (\P) {\txt}; }
     \WedgeArrow{q0}{0}\WedgeArrow{q1}{0}\WedgeArrow{q3}{0}\WedgeArrow{q4}{0}\WedgeArrow{q2}{90}
     \begin{scope}[text=Zcol!85!black, font=\footnotesize]
        \node[below, yshift=-3pt] at (q0) {$Z_0$};
        \node[right, xshift=2pt]  at (q2) {$Z_2$};
     \end{scope}
  \end{scope}

  % Arrow in middle (adjusted for gap)
  \draw[-{Stealth[length=3mm,width=2mm]},red,line width=1.6pt] (-0.6, 0) -- (-0.1, 0);

\end{tikzpicture}
\caption{Logical $Z$}
\end{subfigure}
% END OF PAGE 1
\end{figure}

\begin{figure}[H]\ContinuedFloat
\centering
% ===== MAIN CAPTION HERE =====
\caption{Comparison of 5-qutrit surface code definitions and logical operators. 
(a) The left-hand side (LHS) shows the conventional folded and unfolded definitions from \cite{Moussa2016}. Blue and red regions represent truncated conventional $X$- and $Z$-type stabilizers, respectively. In the folded version, blended colors indicate the mirror duality between these stabilizers. The right-hand side (RHS) illustrates the proposed convention, with the redefined $X$- and $Z$-type stabilizers shown on separate grids using colored patch perimeters. This avoids the folding procedure via local SWAPs, thereby eliminating the need for additional ancillas. 
(b) The proposed logical Hadamard gate, $\bar{H}$. Together with the logical Phase gate $\bar{S}$ shown in (c), these operators generate the single-qutrit logical Clifford group. Unlike the conventional $\bar{H'}$, this implementation requires only nearest-neighbor SWAP gates in Layer 2. 
(c) The proposed logical Phase gate, $\bar{S}$, which utilizes only nearest-neighbor $CZ$ gates, distinct from the conventional $\bar{S'}$. 
(d) Redefined $X$- and $Z$-type stabilizers, selected to be preserved by the logical Clifford group while maintaining mirror dual symmetry. 
(e) The logical $\bar{X}$ operator (unchanged). 
(f) The redefined logical $\bar{Z}$, constructed to commute with stabilizers and generate the Pauli group with $\bar{X}$, ensuring the Pauli group is preserved by the logical Clifford group.}
\label{fig: 5-qutrit-code}

\end{figure}



\subsubsection{Encoding Circuit}

To initialize the system in the code space, we employ a specific unitary circuit that maps the physical ground state $\ket{0}^{\otimes 5}$ to the logical state $\ket{\bar{+}}$. This state corresponds to the $+1$ eigenstate of the logical Pauli operator $\overline{X} = X_0 \otimes X_3$, defined in Eq.~\eqref{eq:logicalPauli}.

\begin{figure}[htbp]
    \centering
    \begin{quantikz}
      \lstick{$q_0: \ket{0}$} & \qw & \qw & \gate{X_3} & \gate{X_3} & \qw \\
      \lstick{$q_1: \ket{0}$} & \qw & \gate{X_3^{\dagger}} & \ctrl{-1} & \qw & \qw \\
      \lstick{$q_2: \ket{0}$} & \gate{H_3^{\dagger}} & \ctrl{-1} & \gate{X_3} & \qw & \qw \\
      \lstick{$q_3: \ket{0}$} & \gate{H_3^{\dagger}} & \qw & \qw & \ctrl{-3}  & \qw \\
      \lstick{$q_4: \ket{0}$} & \gate{H_3^{\dagger}} & \qw & \ctrl{-2}  & \qw  & \qw
    \end{quantikz}
    \caption{Non-fault-tolerant encoding circuit for the $[[5,1,2]]_3$ surface code. The circuit maps the physical state $\ket{0}^{\otimes 5}$ to the logical state $\ket{\bar{+}}$, which is the $+1$ eigenstate of the logical $\bar{X}$ operator. The gates denoted are qutrit Hadamard ($H_3$) and Generalized SUM/CNOT gates. In our simulations this circuit is noiseless (\autoref{sub:noise_models}). On hardware, errors occurring during state preparation contribute SPAM errors to the LRB-D protocol.}
    \label{fig:encoding_circuit}
\end{figure}

The encoding circuit shown in \autoref{fig:encoding_circuit} is a non-fault-tolerant state preparation sequence that minimizes the use of ancillas and the physical gate count. In the simulations of this work the encoding is noiseless, so it contributes no error. On hardware, the quality of the encoding step is classified under State Preparation and Measurement (SPAM) errors. As established in standard RB theory, these SPAM errors are absorbed into the $A$ and $B$ parameters of the fidelity decay curve (Eq.~\eqref{eq:fseq}) and do not affect the extracted decay parameter $f$ or the estimated gate fidelity \cite{Magesan2011PRL,Magesan2012PRA}. With the offset fixed to $B = 0$ (\autoref{sub:fit_model}), they enter through $A$ alone. Consequently, a compact, non-fault-tolerant encoding is sufficient for providing a consistent entry point into the logical code space for benchmarking purposes.

